\section{Bilinear Forms}

\begin{definition}[Bilinear Form]
    Let \(U, V\) be \(\FF\)-vector spaces. A function \(\functiondomain{\phi}{U \times V}{\FF}\) is a \term{bilinear foron} if for all \(\vb{u}_1 \in U\), \(\vb{v}_0 \in V\), the maps
    \[
        \function{\phi_{\vb{u}_0}}{V}{\FF}{\vb{v}}{\phi\pqty{\vb{u}_0, \vb{v}}} \qand \function{\phi_{\vb{v}_0}}{U}{\FF}{\vb{u}}{\phi\pqty{\vb{u}, \vb{v}_0}}
    \]
    are linear maps.
\end{definition}

\begin{examples}
    The following are examples of bilinear forms.

    \begin{enumerate}[label=(\roman*)]
        \item Let \(U = V = \FF^n\), and define
        \[
            \phi\pqty{\vb{u}, \vb{v}} = \sum_{i=1}^n u_i v_i.
        \]

        \item Let \(U = \FF^m\) and \(V = \FF^n\), and \(A \in \matrixset_{m \times n} \pqty{\FF}\). Define
        \[
            \phi\pqty{\vb{u}, \vb{v}} = \vb{u}\transpose A \vb{v}.
        \]

        \item Let \(U = V = \contClass{0} \pqty{\ccintv{0}{1}}\). Define
        \[
            \phi\pqty{f, g} = \int_0^1 f\pqty{x} g\pqty{x} \dd{x}.
        \]

        \item For any \(\FF\)-vector space \(V\), define
        \[
            \function{\phi}{V \times \dual{V}}{\FF}{\phi\pqty{\vb{v}, \theta}}{\theta\pqty{\vb{v}}}.
        \]
    \end{enumerate}
\end{examples}

\subsection{Matrix Representation of Bilinear Forms}

\begin{definition}[Matrix of a Bilinear Form]
    Let \(U, V\) be finite-dimensional vector spaces, with bases \(\basis = \set{\vb{b}_1, \ldots, \vb{v}_m}\) and \(\mathcal{C} = \set{\vb{c}_1, \ldots, \vb{c}_n}\), respectively, then the \term{matrix of \(\phi\) with respect to \(\basis\) and \(\mathcal{C}\)} is the matrix
    \[
        \bqty{\phi}_{\basis, \mathcal{C}} = \pqty{\phi\pqty{\vb{b}_i, \vb{c}_j}}_{i, j} \in \matrixset_{m \times n} \pqty{\FF}.
    \]
\end{definition}

\begin{proposition}[Compatibility of Bilinear Forms and Component Vectors]
    \label{prop:compatibility-bilinear-forms-component-vectors}%
    For all \(\vb{u} \in U\), \(\vb{v} \in V\), we have
    \[
        \bqty{\vb{u}}_{\basis}\transpose \bqty{\phi}_{\basis, \mathcal{C}} \bqty{\vb{v}}_{\mathcal{C}} = \phi\pqty{\vb{u}, \vb{v}}.
    \]

    In particular, the matrix of a bilinear form is unique with respect to the choice of bases.
\end{proposition}

\begin{proof}
    Decompose
    \[
        \vb{u} = \sum_{i=1}^m \lambda_i \vb{b}_i \qand \vb{v} = \sum_{j=1}^n \mu_j \vb{c}_j,
    \]
    so we have \(\bqty{\vb{u}}_{\basis} = \vb*{\lambda}\) and \(\bqty{\vb{v}}_{\mathcal{C}} = \vb*{\mu}\). Then we have
    \begin{align*}
        \phi\pqty{\vb{u}, \vb{v}} &= \phi\pqty{\sum_{i=1}^m \lambda_i \vb{b}_i, \sum_{j=1}^n \mu_j \vb{c}_j} \\
        &= \sum_{i = 1}^{m} \sum_{j = 1}^{n} \lambda_i \mu_j \phi\pqty{\vb{b}_i, \vb{c}_j} \\
        &= \vb*{\lambda}\transpose \bqty{\phi}_{\basis, \mathcal{C}} \vb*{\mu}.
    \end{align*}

    If there is some other matrix \(A \in \matrixset_{m \times n} \pqty{\FF}\) is such that
    \[
        \bqty{\vb{u}}_{\basis}\transpose A \bqty{\vb{v}}_{\mathcal{C}} = \phi\pqty{\vb{u}, \vb{v}},
    \]
    for all \(\vb{u} \in U\), \(\vb{v} \in V\), then by setting \(\vb{u} = \vb{b}_i\) and \(\vb{v} = \vb{c}_j\), then \(\bqty{\vb{u}}_{\basis} = \vb{e}_i\) and \(\bqty{\vb{v}}_{\mathcal{C}} = \vb{e}_j\). So we have
    \[
        \phi\pqty{\vb{b}_i, \vb{c}_j} = \vb{e}_i\transpose A \vb{e}_j = A_{ij},
    \]
\end{proof}

\begin{corollary}[Change of Basis for Bilinear Forms]
    \label{cor:change-of-basis-bilinear-forms}%
    For \(\basis, \basis'\) bases for \(U\), \(\mathcal{C}, \mathcal{C}'\) bases for \(V\), we have
    \[
        \bqty{\phi}_{\basis', \mathcal{C}'} = \pqty{\bqty{\identity_{U}}_{\basis}^{\basis'}}\transpose \bqty{\phi}_{\basis, \mathcal{C}} \bqty{\identity_{V}}_{\mathcal{C}}^{\mathcal{C}'}.
    \]
\end{corollary}

\begin{proof}
    Let \(\vb{u} \in U\), \(\vb{v} \in V\). Then we have
    \begin{align*}
        \bqty{\vb{u}}_{\basis'}\transpose \pqty{\bqty{\identity_{U}}_{\basis}^{\basis'}}\transpose \bqty{\phi}_{\basis, \mathcal{C}} \bqty{\identity_{V}}_{\mathcal{C}}^{\mathcal{C}'} \bqty{\vb{v}}_{\mathcal{C}'} &= \pqty{\bqty{\identity_{U}}_{\basis}^{\basis'} \bqty{\vb{u}}_{\basis'}}\transpose \bqty{\phi}_{\basis, \mathcal{C}} \pqty{\bqty{\identity_{V}}_{\mathcal{C}}^{\mathcal{C}'} \bqty{\vb{v}}_{\mathcal{C}'}} \\
        &= \bqty{\vb{u}}_{\basis}\transpose \bqty{\phi}_{\basis, \mathcal{C}} \bqty{\vb{v}}_{\mathcal{C}} \\
        &= \phi\pqty{\vb{u}, \vb{v}},
    \end{align*}
    hence the results follow by the uniqueness of the matrix.
\end{proof}

\begin{definition}[Rank of a Bilinear Form]
    Let \(\phi\) be a bilinear form on two finite-dimensional vector spaces \(U\) and \(V\). The \term{rank} of a bilinear form \(\phi\) is the rank of its matrix with respect to any choice of bases.
\end{definition}

\begin{remark}
    As a result of \autoref{cor:change-of-basis-bilinear-forms}, since equivalent matrices have the same rank, the rank of a bilinear form is well-defined and independent of the choice of bases.

    For a bilinear form on two identical vector spaces \(U = V\) with bases \(\basis, \basis'\), there is some \(P \in \GL_n\pqty{\FF}\) such that
    \[
        \bqty{\phi}_{\basis', \basis'} = P\transpose \bqty{\phi}_{\basis, \basis} P,
    \]
\end{remark}

\begin{definition}[Congruent Matrices]
    Two matrices \(A, A' \in \matrixset_{n \times n} \pqty{\FF}\) are \term{congruent} if there exists some \(P \in \GL_n\pqty{\FF}\) such that
    \[
        A' = P\transpose A P.
    \]
    
\end{definition}

This is also an equivalence relation on \(\matrixset_{n \times n} \pqty{\FF}\).

\subsection{Left and Right Linear Maps}

\begin{definition}[Left and Right Linear Maps of a Bilinear Form]
    Let \(\functiondomain{\phi}{U \times V}{\FF}\) be a bilinear form. For \(\vb{u} \in U, \vb{v} \in V\), define
    \[
        \phi_{\textL} \pqty{\vb{u}} = \phi\pqty{\vb{u}, \cdot} \in \dual{V} \qand \phi_{\textR}\pqty{\vb{v}} = \phi\pqty{\cdot, \vb{v}} \in \dual{U}.
    \]  

    These are known as the \term{left} and \term{right} linear maps of \(\phi\), respectively.
\end{definition}

\begin{lemma}[Linearity of Left and Right Linear Maps]
    \(\phi_{\text{L}}\) and \(\phi_{\text{R}}\) are linear maps.
\end{lemma}

\begin{definition}[Left and Right Kernels of a Bilinear Form]
    Let \(\functiondomain{\phi}{U \times V}{\FF}\) be a bilinear form. The \term{left kernel} of \(\phi\) is
    \[
        \ker_{\textL} \phi = \set{\vb{u} \in U}{\phi_{\text{L}}\pqty{\vb{u}} = 0} = \set{\vb{u} \in U }{\forall \vb{v} \in V: \phi\pqty{\vb{u}, \vb{v}} = 0},
    \]
    and the \term{right kernel} of \(\phi\) is
    \[
        \ker_{\textR} \phi = \set{\vb{v} \in V}{\phi_{\text{R}}\pqty{\vb{v}} = 0} = \set{\vb{v} \in V}{\forall \vb{u} \in U : \phi\pqty{\vb{u}, \vb{v}} = 0}.
    \]
\end{definition}

\begin{proposition}[Matrix Representation of Left and Right Linear Maps]
    \label{prop:matrix-representation-left-right-linear-maps}%
    Let \(U, V\) be finite-dimensional vector spaces, and let \(\basis, \mathcal{C}\) be bases for \(U\) and \(V\), respectively. Let \(\dual{\basis}\) and \(\dual{\mathcal{C}}\) be the dual bases for \(\dual{U}\) and \(\dual{V}\), respectively. Then we have
    \[
        \bqty{\phi_{\textL}}^{\basis}_{\dual{\mathcal{C}}} = \bqty{\phi}_{\basis, \mathcal{C}}\transpose \qand \bqty{\phi_{\textR}}^{\mathcal{C}}_{\dual{\basis}} = \bqty{\phi}_{\basis, \mathcal{C}}.
    \]
\end{proposition}

\begin{proof}
    Let \(A = \bqty{\phi}_{\basis, \mathcal{C}}\), and \(A' = \bqty{\phi_{\textL}}^{\basis}_{\dual{\mathcal{C}}}\). Let \(\basis = \set{\vb{b}_1, \ldots, \vb{b}_m}\) and \(\mathcal{C} = \set{\vb{c}_1, \ldots, \vb{c}_n}\). Then we have
    \[
        \phi_{\textL}\pqty{\vb{b}_i} = \sum_{k = 1}^{n} a'_{kj} \dual{\vb{c}}_k,
    \]
    while we have
    \begin{align*}
        a_{ji} &= \phi \pqty{\vb{b}_j, \vb{c}_i} \\
        &= \phi_{\textL}\pqty{\vb{b}_j}\pqty{\vb{c}_i} \\
        &= \sum_{k = 1}^{n} a'_{kj} \dual{\vb{c}}_k\pqty{\vb{c}_i} \\
        &= a'_{ij},
    \end{align*}
    so we have \(A' = A\transpose\). The proof for the right linear map is similar.
\end{proof}

\begin{corollary}[Rank of Left and Right Linear Maps]
    \label{cor:rank-of-left-and-right-linear-maps}%
    Let \(\functiondomain{\phi}{U \times V}{\FF}\) be a bilinear form. Then we have
    \[
        \rank \phi_{\textL} = \rank \phi_{\textR} = \rank \phi.
    \]
\end{corollary}

\begin{proof}
    This is immediate from \autoref{prop:matrix-representation-left-right-linear-maps} and \autoref{thm:row-rank-equals-column-rank} (Row Rank equals Column Rank).
\end{proof}

\begin{definition}[Orthogonal Complement/Perpendicular Subspace]
    Let \(\functiondomain{\phi}{U \times V}{\FF}\) be a bilinear form. Let \(S \subseteq U\) and \(T \subseteq V\). The \term{orthogonal complement/perpenducular subspace} of \(S\) is
    \[
        \rperp{S} = \set{\vb{v} \in V}{\forall \vb{u} \in S: \phi\pqty{\vb{u}, \vb{v}} = 0} \subseteq V,
    \]
    and of \(T\) is
    \[
        \lperp{T} = \set{\vb{u} \in U}{\forall \vb{v} \in T: \phi\pqty{\vb{u}, \vb{v}} = 0} \subseteq U.
    \]
\end{definition}

\begin{remarks}
    \begin{itemize}
        \item By the subspace test, \(\rperp{S} \subspace V\) and \(\lperp{T} \subspace U\).
        \item If \(S_1 \subseteq S_2 \subseteq U\), then \(\rperp{S_2} \subseteq \rperp{S_1}\). Similarly, if \(T_1 \subseteq T_2 \subseteq V\), then \(\lperp{T_2} \subseteq \lperp{T_1}\).
        \item \(\ker \phi_{\textL} = \lperp{V}\) and \(\ker \phi_{\textR} = \rperp{U}\).
    \end{itemize}
\end{remarks}

\begin{definition}[Degeneracy of Bilinear Forms]
    A bilinear form \(\functiondomain{\phi}{U \times V}{\FF}\) is \term{degenerate} if either \(\rperp{U}\) or \(\lperp{V}\) is non-trivial. Otherwise, it is \term{non-degenerate}.
\end{definition}

\begin{proposition}[Non-Degeneracy of Bilinear Forms]
    For finite-dimensional vector spaces \(U, V\) with bases \(\basis\) and \(\mathcal{C}\), a bilinear form \(\functiondomain{\phi}{U \times V}{\FF}\) is non-degenerate if and only if \(\dim U = \dim V\), and the matrix \(\phi_{\basis, \mathcal{C}}\) is invertible.
\end{proposition}

\begin{proof}
    Non-degeneracy is equivalent to \(\ker \phi_{\textL} = \set{\vb{0}}\) and \(\ker \phi_{\textR} = \set{\vb{0}}\). By \autoref{thm:rank-nullity-thm} (Rank-Nullity Theorem), we have
    \[
        \dim U = \rank \phi_{\textL} = \rank \phi = \rank \phi_{\textR} = \dim V,
    \]
    by \autoref{cor:rank-of-left-and-right-linear-maps}. This happens if and only if \(\phi_{\basis, \mathcal{C}}\) is a square matrix of full rank.
\end{proof}

\subsection{Bilinear Form on a Single Vector Space}

If \(\functiondomain{\phi}{U \times V}{\FF}\) is a bilinear form where \(U = V\) are finite-dimensional with basis \(\basis\), then we denote
\[
    \bqty{\phi}_{\basis} = \bqty{\phi}_{\basis, \basis} \in \matrixset_{n \times n} \pqty{\FF}.
\]

\begin{definition}[Symmetric Bilinear Form]
    A bilinear form \(\functiondomain{\phi}{V \times V}{\FF}\) is \term{symmetric} if \(\phi\pqty{\vb{u}, \vb{v}} = \phi\pqty{\vb{v}, \vb{u}}\) for all \(\vb{u}, \vb{v} \in V\).
\end{definition}

\begin{remarks}
    \begin{itemize}
        \item For a finite-dimensional vector space \(V\) with basis \(\basis\), a bilinear form \(\functiondomain{\phi}{V \times V}{\FF}\) is symmetric if and only if \(\bqty{\phi}_{\basis}\) is a symmetric matrix.
        \item If \(\phi\) is symmetric, then for any \(S \subseteq V\), we have \(\lperp{S} = \rperp{S}\).
        \item In particular, if \(\phi\) is symmetric, then \(\ker \phi_{\textL} = \ker \phi_{\textR}\). We call this the \term{kernel} of \(\phi\).
    \end{itemize}
\end{remarks}

One cool result about symmetric bilinear forms is that we can recover a symmetric bilinear form from only looking at the values \(\phi\pqty{\vb{v}, \vb{v}}\) for all \(\vb{v} \in V\). This motivates the following definition.

\begin{definition}[Quadratic Form]
    A function \(\functiondomain{Q}{V}{\FF}\) is a \term{quadratic form} if there exists a symmetric bilinear form \(\functiondomain{\phi}{V \times V}{\FF}\) such that for all \(\vb{v} \in V\), we have \(Q\pqty{\vb{v}} = \phi\pqty{\vb{v}, \vb{v}}\).
\end{definition}

\begin{example}
    Let \(V = \RR^2\), for \(A = \pmqty{a & b \\ c & d} \in \matrixset_{2 \times 2} \pqty{\RR}\), define the quadratic form
    \[
        Q_A \pqty{\vb{v}} = \vb{v}\transpose A \vb{v} = a v_1^2 + \pqty{b + c} v_1 v_2 + d v_2^2.
    \]

    There is a corresponding bilinear form
    \[
        \phi_A \pqty{\vb{u}, \vb{v}} = \vb{u}\transpose A \vb{v}.
    \]

    Note that \(Q_A = Q_{\frac{1}{2} \pqty{A + A\transpose}}\), so we can always assume that \(A\) is symmetric. (The asymmetric part of \(A\) does not contribute to the quadratic form.)
\end{example}

\begin{proposition}[Polarisation Identity]
    \label{prop:polarisation-identity}%
    Suppose \(1 + 1 \neq 0\) in \(\FF\), and let \(V\) be a \(\FF\)-vector space. Let \(\functiondomain{Q}{V}{\FF}\) be a quadratic form. Then there is a unique symmetric bilinear form \(\functiondomain{\psi}{V \times V}{\FF}\) such that for all \(\vb{v} \in V\), we have \(Q\pqty{\vb{v}} = \psi\pqty{\vb{v}, \vb{v}}\), which is given by
    \[
        \psi\pqty{\vb{u}, \vb{v}} = \frac{1}{2} \pqty{Q\pqty{\vb{u} + \vb{v}} - Q\pqty{\vb{u}} - Q\pqty{\vb{v}}}.
    \]
\end{proposition}

\begin{proof}
    To show existence, let \(\phi\) induce this quadratic form \(Q\), i.e., \(\phi\pqty{\vb{v}, \vb{v}} = Q\pqty{\vb{v}}\) for all \(\vb{v} \in V\). Set \(\psi\pqty{\vb{u}, \vb{v}} = \frac{1}{2} \pqty{\phi\pqty{\vb{u}, \vb{v}} + \phi\pqty{\vb{v}, \vb{u}}}\). Then \(\psi\) is symmetric, and for all \(\vb{v} \in V\), we have \(\psi\pqty{\vb{v}, \vb{v}} = \phi\pqty{\vb{v}, \vb{v}} = Q\pqty{\vb{v}}\).

    For uniqueness, suppose \(\psi\) is a symmetric bilinear form satisfying \(\psi\pqty{\vb{v}, \vb{v}} = Q\pqty{\vb{v}}\) for all \(\vb{v} \in V\). Then note 
    \begin{align*}
        Q\pqty{\vb{u} + \vb{v}} &= \psi\pqty{\vb{u} + \vb{v}, \vb{u} + \vb{v}} \\
        &= \psi\pqty{\vb{u}, \vb{u}} + \psi\pqty{\vb{u}, \vb{v}} + \psi\pqty{\vb{v}, \vb{u}} + \psi\pqty{\vb{v}, \vb{v}} \\
        &= Q\pqty{\vb{u}} + \psi\pqty{\vb{u}, \vb{v}} + \psi\pqty{\vb{v}, \vb{u}} + Q\pqty{\vb{v}} \\
        &= Q\pqty{\vb{u}} + 2 \psi\pqty{\vb{u}, \vb{v}} + Q\pqty{\vb{v}},
    \end{align*}
    and rearranging, we get that \(\psi\) is uniquely determined by \(q\) as
    \[
        \psi\pqty{\vb{u}, \vb{v}} = \frac{1}{2} \pqty{Q\pqty{\vb{u} + \vb{v}} - Q\pqty{\vb{u}} - Q\pqty{\vb{v}}},
    \]
\end{proof}

\begin{theorem}[Diagonalisation of Symmetric Bilinear Forms]
    \label{thm:diagonalisation-symmetric-bilinear-forms}%
    Suppose \(1 + 1 \neq 0\) in \(\FF\). Let \(V\) be a finite-dimensional \(\FF\)-vector space, and let \(\functiondomain{\phi}{V \times V}{\FF}\) be a symmetric bilinear form. Then there exists a basis \(\basis\) for \(V\) such that \(\bqty{\phi}_{\basis}\) is diagonal.
\end{theorem}

\begin{remark}
    The assumption that \(\phi\) is symmetric is necessary, since if \(A \in \matrixset_{n \times n} \pqty{\FF}\) with \(P \transpose A P = D\) being diagonal for some \(P \in \GL_n\pqty{\FF}\), then we have
    \[
        P\transpose A P = D = D\transpose = P\transpose A\transpose P
    \]
    so \(A = A\transpose\), i.e., \(A\) is symmetric.
\end{remark}

\begin{corollary}[Congruence of Symmetric Matrices to Diagnoal Matrices]
    Suppose \(1 + 1 \neq 0\) in \(\FF\), every symmetric matrix over \(\FF\) is congruent to a diagonal matrix.
\end{corollary}

\begin{example}
    If \(\FF = \flatfrac{\ZZ}{2 \ZZ}\), the matrix \(\pmqty{0 & 1 \\ 1 & 0}\) is not congruent to any diagonal matrix.
\end{example}

\begin{proof}[\proofname\ of \autoref{thm:diagonalisation-symmetric-bilinear-forms}]
    We induct on \(n = \dim V\). This is clear for \(n = 0\) or \(1\).

    If \(\phi\pqty{\vb{v}, \vb{v}} = 0\) for all \(\vb{v} \in V\), then by \autoref{prop:polarisation-identity} (Polarisation Identity), we have \(\phi\pqty{\vb{u}, \vb{v}} = 0\) for all \(\vb{u}, \vb{v} \in V\), so any basis will do.

    Now, suppose \(\vb{v}_1 \in V\) with \(\phi\pqty{\vb{v}_1, \vb{v}_1} \neq 0\). Let
    \[
        U = \rperp{\set{\vb{v}_1}} = \ker \pqty{\phi_{\textL}\pqty{\vb{v}_1}} \subseteq V.
    \]

    Note \(\vb{v}_1 \notin U\), so by \autoref{thm:rank-nullity-thm} (Rank--Nullity Theorem), we have \(\dim U = n - 1\), and \(V = \spn\set{\vb{v}_1} \oplus U\).

    By the induction hypothesis, there is a basis \(\basis_U = \set{\vb{v}_2, \vb{v}_2, \ldots, \vb{v}_{n}}\) for \(U\), such that the restriction of \(\phi\) to \(U\) is diagonal with respect to \(\basis_U\).

    Set \(\basis = \set{\vb{v}_1} \cup \basis_U\). Then \(\basis\) is a basis for \(V\), and we have
    \[
        \bqty{\phi}_{\basis} = \pmqty{\phi\pqty{\vb{v}_1, \vb{v}_1} & \vb{0}\transpose \\ \vb{0} & \bqty{\phi}_{\basis_U}},
    \]
    is indeed diagonal, by the construction of the orthogonal complement \(U\) of \(\vb{v}_1\).
\end{proof}

\begin{remark}
    This is a procedure that will always work for constructing a diagonal basis. In practice, it is often easier to do by inspection.
\end{remark}

\begin{example}
    For \(\vb{x} \in \RR^3\), define the quadratic form
    \begin{align*}
        Q\pqty{\vb{x}} &= x_1^2 - 4 x_2 x_3 + x_3^2 \\
        &= x_1^2 + \pqty{x_3 - 2 x_2}^2 - \pqty{2 x_2}^2 \\
        &= z_1^2 + z_2^2 - z_3^2,
    \end{align*}
    where
    \[
        \vb{z} = \pmqty{1 & 0 & 0 \\ 0 & -2 & 1 \\ 0 & 2 & 0} \vb{x} = P\inverse \vb{x}.
    \]

    Then \(Q \pqty{\vb{x}} = \psi\pqty{\vb{x}, \vb{x}}\), where
    \[
        \psi\pqty{\vb{x}, \vb{y}} = x_1 y_1 - 2 \pqty{x_2 y_3 + x_3 y_2} + x_3 y_3.
    \]

    In the standard basis, \(\psi\) has matrix
    \[
        A = \pmqty{1 & 0 & 0 \\ 0 & 0 & -2 \\ 0 & -2 & 1},
    \]
    then we have
    \[
        P\transpose A P = \pmqty{1 & 0 & 0 \\ 0 & 1 & 0 \\ 0 & 0 & -1},
    \]
    is diagonal as desired.
\end{example}

\begin{corollary}[Diagonalisation of Symmetric Bilinear Forms over \(\RR\) and \(\CC\)]
    \label{cor:diagonalisation-symmetric-bilinear-forms-over-r-and-c}%
    Suppose \(\FF = \CC\) or \(\RR\), and let \(U, V\) and \(\phi\). be as in \autoref{thm:diagonalisation-symmetric-bilinear-forms}. Define \(Q\pqty{\vb{v}} = \phi\pqty{\vb{v}, \vb{v}}\) be the quadratic form induced by \(\phi\).

    If \(\FF = \CC\), then there exists a basis \(\basis = \set{\vb{b}_1, \vb{b}_2, \ldots, \vb{b}_n}\) for \(V\) such that
    \[
        \bqty{\phi}_{\basis} = \pmqty{\dmat{I_r, O}} \qfor r = \rank \phi,
    \]
    and equivalently
    \[
        Q\pqty{\sum_{i = 1}^{n} \lambda_i \vb{v}_i} = \sum_{i = 1}^{r} \lambda_i^2.
    \]

    If \(\FF = \RR\), then there exists a basis \(\basis\) for \(V\) such that
    \[
        \bqty{\phi}_{\basis} = \pmqty{\dmat{I_p, -I_q, O}} \qfor p + q = \rank \phi,
    \]
    and equivalently
    \[
        Q\pqty{\sum_{i = 1}^{n} \lambda_i \vb{v}_i} = \sum_{i = 1}^{p} \lambda_i^2 - \sum_{i = p + 1}^{p + q} \lambda_i^2.
    \]
\end{corollary}

\begin{proof}
    Pick bases \(\mathcal{C} = \set{\vb{w}_1, \ldots, \vb{w}_n}\) as in \autoref{thm:diagonalisation-symmetric-bilinear-forms}. Them we have
    \[
        \bqty{\phi}_{\mathcal{C}} = \pmqty{\dmat{a_1, \ddots, a_n}},
    \]
    then by reordering \(\mathcal{C}\), we may assume that:
    \item \(a_i \neq 0\) for \(i \leq r = \rank \phi\); and
    \item in the case \(\FF = \RR\), \(r = p + q\), and \(a_i > 0\) for \(i \leq p\), and \(a_i < 0\) for \(p < i \leq p + q\).
    
    If \(\FF = \CC\), rescale the vectors, such that
    \[
        \vb{v}_i = \begin{cases}
            \flatfrac{\vb{w}_i}{\sqrt{a_i}}, & i \leq r, \\
            \vb{w}_i, & i > r,
        \end{cases}
    \]
    and if \(\FF = \RR\), rescale the vectors, such that
    \[
        \vb{v}_i = \begin{cases}
            \flatfrac{\vb{w}_i}{\sqrt{a_i}}, & i \leq p, \\
            \flatfrac{\vb{w}_i}{\sqrt{-a_i}}, & p < i \leq p + q, \\
            \vb{w}_i, & i > p + q.
        \end{cases}
    \]

    Then under this basis, the matrix is in the form, as claimed.
\end{proof}

\subsection{Bilinear Forms on \(\RR\)}

\begin{definition}[Positive/Negative-(Semi)Definite Bilinear Forms]
    Let \(V\) be a real vector space, and let \(\functiondomain{\phi}{V \times V}{\RR}\) be a symmetric bilinear form. We say:
    \begin{itemize}
        \item \(\phi\) is \term{positive definite} if its associated quadratic form satisfies \(\phi\pqty{\vb{v}, \vb{v}} > 0\) for all \(\vb{v} \in V \setminus \set{\vb{0}}\).
        \item \(\phi\) is \term{positive semidefinite} if its associated quadratic form satisfies \(\phi\pqty{\vb{v}, \vb{v}} \geq 0\) for all \(\vb{v} \in V\).
        \item \(\phi\) is \term{negative definite} if its associated quadratic form satisfies \(\phi\pqty{\vb{v}, \vb{v}} < 0\) for all \(\vb{v} \in V \setminus \set{\vb{0}}\).
        \item \(\phi\) is \term{negative semidefinite} if its associated quadratic form satisfies \(\phi\pqty{\vb{v}, \vb{v}} \leq 0\) for all \(\vb{v} \in V\).
    \end{itemize}

    Otherwise, we say \(\phi\) is \term{indefinite}.
\end{definition}

\begin{example}
    Let \(V\) be a finite-dimensional real vector space, and let \(\functiondomain{\phi}{V \times V}{\RR}\) be a symmetric bilinear form. Then
    \begin{enumerate}[label=(\roman*)]
        \item If
        \[
            \bqty{\phi}_{\basis} = \pmqty{\dmat{I_p, O}},
        \]
        the \(\phi\) is positive semidefinite, and it is positive definite if and only if \(p = \dim V\).
        \item In the general case, if
        \[
            \bqty{\phi}_{\basis} = \pmqty{\dmat{I_p, -I_q, O}},
        \]
        where \(p, q > 0\), we have that \(\phi\) is indefinite.
    \end{enumerate}
\end{example}

\begin{theorem}[Sylvester's Law of Inertia]
    \label{thm:sylvesters-law-of-inertia}%
    Let \(V\) be a finite-dimensional real vector space, and let \(\functiondomain{\phi}{V \times V}{\RR}\) be a symmetric bilinear form. If \(\basis, \basis'\) are bases for \(V\) such that
    \[
        \bqty{\phi}_{\basis} = \pmqty{\dmat{I_p, -I_q, O}} \qand \bqty{\phi}_{\basis'} = \pmqty{\dmat{I_{p'}, -I_{q'}, O}},
    \]
    then \(p = p'\) and \(q = q'\).
\end{theorem}

\begin{definition}[Signature of a Bilinear Form]
    Let \(V\) be a finite-dimensional real vector space, and let \(\functiondomain{\phi}{V \times V}{\RR}\) be a symmetric bilinear form. The \term{signature} of \(\phi\) is \(\sigma\pqty{\phi} = p - q\), where \(p\) and \(q\) are as in Sylvester's Law of Inertia.
\end{definition}

\begin{corollary}[Well-Definedness of Signature]
    The signature of a symmetric bilinear form is well-defined and independent of the choice of basis.
\end{corollary}

To prove Sylvester's law, we introduce the following lemma.

\begin{lemma}[Trivial Intersection of Subspaces with Positive and Negative Definite Restrictions]
    Let \(U, W \subspace V\) where \(V\) is a real vector space, and let \(\functiondomain{\phi}{V \times V}{\RR}\) be a symmetric bilinear form.
    \begin{enumerate}[label=(\roman*)]
        \item If the restriction \(\functiondomain{\rightbar{\phi}_{U}}{U \times U}{\RR}\) is positive definite, and \(\functiondomain{\rightbar{\phi}_{W}}{W \times W}{\RR}\) is negative semidefinite, then \(U \cap W = \set{\vb{0}}\) intersect trivially.
        \item If the restriction \(\functiondomain{\rightbar{\phi}_{U}}{U \times U}{\RR}\) is negative definite, and \(\functiondomain{\rightbar{\phi}_{W}}{W \times W}{\RR}\) is positive semidefinite, then \(U \cap W = \set{\vb{0}}\) intersect trivially.
    \end{enumerate}
\end{lemma}

\begin{proof}
    \begin{enumerate}[label=(\roman*)]
        \item Suppose \(\vb{0} \neq \vb{v} \in U \cap W\). Then we have
        \[
            \vb{v} \in U \implies \phi\pqty{\vb{v}, \vb{v}} > 0 \qand \vb{v} \in W \implies \phi\pqty{\vb{v}, \vb{v}} \leq 0,
        \]
        which is a contradiction. Hence \(U \cap W = \set{\vb{0}}\).
        \item This is similar to (i).
    \end{enumerate}
\end{proof}


\begin{proof}[\proofname\ of \autoref{thm:sylvesters-law-of-inertia} (Sylvester's Law)]
    For uniqueness of \(p\), we claim that \(p\) is the largest dimension of a subspace of \(V\), on which \(\phi\) is positive definite.

    Set \(U = \spn \set{\vb{v}_1, \ldots, \vb{v}_p}\), and \(W = \set{\vb{v}_{p + 1}, \ldots, \vb{v}_n}\). Then \(\rightbar{\phi}_U\) is positive definite, while \(\rightbar{\phi}_W\) is negative semidefinite.

    Suppose there is another subspace \(U' \subspace V\) with \(\rightbar{\phi}_{U'}\) being positive definite, but \(\dim U' = p' > p\). Then by the previous lemma, we have \(U \cap W = \set{\vb{0}}\), but by the dimension formula, we have
    \[
        \dim \pqty{U'} + \dim \pqty{W} > n = \dim V,
    \]
    so \(U' \cap W \neq \set{\vb{0}}\), which is a contradiction. Hence \(p' \leq p\), and by symmetry, we have \(p = p'\).

    The uniqueness of \(q\) is similar.
\end{proof}

\begin{definition}[Totally Isotropic Subspace]
    Let \(V\) be a finite-dimensional real vector space, and let \(\functiondomain{\phi}{V \times V}{\RR}\) be a symmetric bilinear form. A subspace \(T \subspace V\) is \term{totally isotropic} if the restriction \(\rightbar{\phi}_T\) is identically zero, i.e., \(\phi\pqty{\vb{t}_1, \vb{t}_2} = 0\) for all \(\vb{t}_1, \vb{t}_2 \in T\).
\end{definition}

\begin{example}
    Suppose \(V\) is finite-dimensional with \(\basis\), and
    \[
        \bqty{\phi}_{\basis} = \pmqty{\dmat{I_p, -I_q, O}}.
    \]

    Then the maximal dimension of the totally isotropic subspace is not only \(\dim{V} - p - q\), but in fact is \(\dim \pqty{V} - \max\pqty{p, q}\).
\end{example}

\subsection{Sesquilinear Forms on \(\CC\)}

The motivating example for bilinear forms is the dot product on \(\RR^n\). However, the dot product on \(\CC^n\) is not bilinear, but rather a sesquilinear form (recall we take the complex conjugate of one argument).

\begin{example}[Complex Dot Product]
    The standard dot product on \(\CC^n\)
    \[
        \vb{x} \vdot \vb{y} = \sum_{i = 1}^{n} x_i \conjugate{y_i}.
    \]

    (We use the standard algebraic convention of taking the complex conjugate of the second argument.) This is clearly not bilinear.
\end{example}

\begin{definition}[Sesquilinear Form]
    Let \(V, W\) be complex vector spaces. A \term{sesquilinear form} is a function
    \[
        \functiondomain{\phi}{V \times W}{\CC}
    \]
    such that \(\phi\) is linear in the first argument, and conjugate-linear in the second argument, i.e., for all \(\vb{v}_1, \vb{v}_2 \in V\), \(\vb{w}_1, \vb{w}_2 \in W\), and \(\lambda, \mu \in \CC\), we have
    \begin{align*}
        \phi\pqty{\lambda \vb{v}_1 + \vb{v}_2, \vb{w}_1} = \lambda \phi\pqty{\vb{v}_1, \vb{w}_1} + \phi\pqty{\vb{v}_2, \vb{w}_1}, \\
        \phi\pqty{\vb{v}_1, \mu \vb{w}_1 + \vb{w}_2} = \conjugate{\mu} \phi\pqty{\vb{v}_1, \vb{w}_1} + \phi\pqty{\vb{v}_1, \vb{w}_2}.
    \end{align*}
\end{definition}

\begin{definition}[Matrix Representation of Sesquilinear Forms]
    Let \(V, W\) be finite-dimensional complex vector spaces, and let \(\functiondomain{\phi}{V \times W}{\CC}\) be a sesquilinear form. Let \(\basis\) and \(\mathcal{C}\) be bases for \(V\) and \(W\), respectively. The \term{matrix representation} of \(\phi\) with respect to the bases \(\basis\) and \(\mathcal{C}\) is
    \[
        \bqty{\phi}_{\basis, \mathcal{C}} = \pmqty{\phi\pqty{\vb{b}_i, \vb{c}_j}}_{1 \leq i \leq m, 1 \leq j \leq n} \in \matrixset_{m \times n} \pqty{\CC},
    \]
\end{definition}

\begin{proposition}[Compatibility of Sesquilinear Forms and Component Vectors]
    Let \(\basis, \mathcal{C}\) be bases for finite-dimensional complex vector spaces \(V\) and \(W\), respectively. The matrix \(\bqty{\phi}_{\basis, \mathcal{C}}\) is compatible with the component vectors of \(\vb{v} \in V\) and \(\vb{w} \in W\) in the sense that it is the unique matrix such that
    \[
        \bqty{\vb{v}}_{\basis}\transpose \bqty{\phi}_{\basis, \mathcal{C}} \conjugate{\bqty{\vb{w}}_{\mathcal{C}}} = \phi\pqty{\vb{v}, \vb{w}},
    \]
    for all \(\vb{v} \in V\) and \(\vb{w} \in W\).
\end{proposition}

\begin{proof}
    This is similar to the proof of \autoref{prop:compatibility-bilinear-forms-component-vectors}.
\end{proof}

\begin{corollary}[Change of Basis for Sesquilinear Forms]
    Let \(\basis, \basis'\) be bases for a finite-dimensional complex vector space \(V\), and let \(\mathcal{C}, \mathcal{C}'\) be bases for a finite-dimensional complex vector space \(W\). Then we have
    \[
        \bqty{\phi}_{\basis', \mathcal{C}'} = P\transpose \bqty{\phi}_{\basis, \mathcal{C}} \conjugate{Q},
    \]
    where \(P = \bqty{\identity_V}_{\basis}^{\basis'}\) and \(Q = \bqty{\identity_W}_{\mathcal{C}}^{\mathcal{C}'}\).
\end{corollary}

\begin{proof}
    This is similar to the proof of \autoref{cor:change-of-basis-bilinear-forms}.
\end{proof}

\begin{definition}[Hermitian Sesquilinear Forms]
    A sesquilinear form \(\functiondomain{\phi}{V \times V}{\CC}\) is \term{Hermitian} if \(\phi\pqty{\vb{v}, \vb{w}} = \conjugate{\phi\pqty{\vb{w}, \vb{v}}}\) for all \(\vb{v}, \vb{w} \in V\).
\end{definition}

\begin{remarks}
    Let \(V\) be a finite-dimensional complex vector space with basis \(\basis\), and let \(\functiondomain{\phi}{V \times V}{\CC}\) be Hermitian. Then
    \begin{enumerate}[label=(\roman*)]
        \item We denote \(\bqty{\phi}_{\basis} = \bqty{\phi}_{\basis, \basis}\).
        \item For all \(\vb{v} \in V\), we have \(\phi\pqty{\vb{v}, \vb{v}} \in \RR\), and for \(\lambda \in \CC\), \(\phi\pqty{\lambda \vb{v}, \lambda \vb{v}} = \abs{\lambda}^2 \phi\pqty{\vb{v}, \vb{v}}\). This allows us to define positive/negative-(semi)definite Hermitian sesquilinear forms, in the same way as we did for symmetric bilinear forms on \(\RR\).
        \item Note \(\phi\pqty{\vb{u}, \vb{v}} = 0\) if and only if \(\phi\pqty{\vb{v}, \vb{u}} = 0\). Hence, the left and right orthogonal complements of a subspace \(S \subspace V\) are equal.
    \end{enumerate}
\end{remarks}

\begin{definition}[Hermitian Matrix]
    A square complex matrix \(A \in \matrixset_{n \times n} \pqty{\CC}\) is \term{Hermitian} if \(A = A\hermitian\), i.e., \(A = \conjugate{A\transpose}\).
\end{definition}

\begin{lemma}[Hermitian Sesquilinear Forms and Hermitian Matrices]
    \label{lem:hermitian-sesquilinear-forms-hermitian-matrices}%
    Let \(V\) be a finite-dimensional complex vector space with basis \(\basis\), and let \(\functiondomain{\phi}{V \times V}{\CC}\) be a sesquilinear form. Then \(\phi\) is Hermitian, if and only if \(\bqty{\phi}_{\basis}\) is a Hermitian matrix for some basis (equivalently, all bases).
\end{lemma}

\begin{proof}
    Let \(\basis = \set{\vb{b}_1, \ldots, \vb{b}_n}\) be a basis for \(V\), and \(A = \bqty{\phi}_{\basis}\).

    Suppose \(\phi\) is Hermitian. Then we have \(a_{ij} = \phi\pqty{\vb{b}_i, \vb{b}_j} = \conjugate{\phi\pqty{\vb{b}_j, \vb{b}_i}} = \conjugate{a_{ji}}\), so \(A = A\hermitian\).

    Suppose \(A\) is Hermitian. Then for all \(\vb{v}, \vb{w} \in V\), we have
    \begin{align*}
        \phi\pqty{\vb{v}, \vb{w}} &= \bqty{\vb{v}}_{\basis} \transpose A \conjugate{\bqty{\vb{w}}_{\basis}} \\
        &= \pqty{\bqty{\vb{v}}_{\basis} \transpose A \conjugate{\bqty{\vb{w}}_{\basis}}}\transpose \\
        &= \conjugate{\bqty{\vb{w}}_{\basis}} \transpose A\transpose \bqty{\vb{v}}_{\basis} \\
        &= \conjugate{\bqty{\vb{w}}_{\basis}\transpose \conjugate{A}\transpose \conjugate{\bqty{\vb{v}}_{\basis}}} \\
        &= \conjugate{\phi\pqty{\vb{w}, \vb{v}}},
    \end{align*}
    since \(A = A\hermitian = \conjugate{A}\transpose\). Hence \(\phi\) is Hermitian.
\end{proof}

\begin{proposition}[Polarisation Identity for Hermitian Forms]
    \label{prop:polarisation-identity-hermitian-sesquilinear-forms}%
    Let \(V\) be a complex vector space. A Hermitian form \(\phi\) on \(V\) is uniquely determined by its associated quadratic form \(Q\pqty{\vb{v}} = \phi\pqty{\vb{v}, \vb{v}}\), and is given by
    \[
        \phi\pqty{\vb{u}, \vb{v}} = \frac{1}{4} \pqty{Q\pqty{\vb{u} + \vb{v}} - Q\pqty{\vb{u} - \vb{v}} + i Q\pqty{\vb{u} + i \vb{v}} - i Q\pqty{\vb{u} - i \vb{v}}}.
    \]
\end{proposition}

\begin{proof}
    \textbf{\textcolor{red}{See Example Sheet.}}
\end{proof}

\begin{theorem}[Diagonalisation of Hermitian Forms]
    Let \(V\) be a finite-dimensional complex vector space, and let \(\phi\) be a /Hermitian form on \(V\).

    \begin{enumerate}[label=(\roman*)]
        \item There is a basis \(\basis\) of \(V\), such that
        \[
            \bqty{\phi}_{\basis} = \pmqty{\dmat{I_p, -I_q, O}} \qfor p + q = \rank \phi.
        \]

        \item \term{Hermitian Sylvester's Law.} The numbers \(p\) and \(q\) are uniquely determined by \(\phi\), and are independent of the choice of basis.
    \end{enumerate}
\end{theorem}

\begin{proof}
    This is almost identical to the proof fo bilinear forms.

    \textbf{Existence.} (Similar to \autoref{cor:diagonalisation-symmetric-bilinear-forms-over-r-and-c}.) We induct on \(n \dim V\). If \(\phi\) is identically zero, then we are done. If not, then there is some non-zero \(\vb{v}_1 \in V\) such that \(\phi\pqty{\vb{v}_1, \vb{v}_1} \neq 0\), by \autoref{prop:polarisation-identity-hermitian-sesquilinear-forms} (Polarisation Identity for Hermitian Forms). Let \(U = \rperp{\spn\pqty{\vb{v}_1}}\). We may check that \(V = \spn\pqty{\vb{v}_1} \oplus U\), and by the induction hypothesis on \(\rightbar{\phi}_{U}\), we can find a basis for \(U\) such that the restriction of \(\phi\) to \(U\) is diagonal. Then, we rescale basis vectors using the observation that \(\phi\pqty{\lambda \vb{v}, \lambda \vb{v}} = \abs{\lambda}^2 \phi\pqty{\vb{v}, \vb{v}}\), to get the desired diagonal form.

    \textbf{Uniqueness.} (Similar to \autoref{thm:sylvesters-law-of-inertia}.) The uniqueness of \(p\) and \(q\) follows from the fact that \(p\) is the largest dimension of a subspace of \(V\) on which \(\phi\) is positive definite, and \(q\) is the largest dimension of a subspace of \(V\) on which \(\phi\) is negative definite.
\end{proof}